%%%%%%%%%%%%%%%%%%%%%%%%%%%%%%%%%%%%%%%%%%%%%%%%%%%%%%%%%%%%%%%%%%%%%%
%% Shared pieces for the l04_afe signal flow graphs, l4_first_order and
%% l4_biquad.  The biquad is meant to read as the first-order graph with a
%% second integrator added, so the summing node and the 1/s block have to be
%% the same size and shape in both.  They are the same size in the original
%% artwork too.
%%
%% Names are prefixed and letters only, and computed coordinates are braced:
%% TikZ scans a coordinate for its closing parenthesis, so a '(' inside an
%% expression ends it early.
%%%%%%%%%%%%%%%%%%%%%%%%%%%%%%%%%%%%%%%%%%%%%%%%%%%%%%%%%%%%%%%%%%%%%%

\newcommand{\sfgR}{0.25}    % summing node radius
\newcommand{\sfgBw}{1.0}    % half width of the 1/s block
\newcommand{\sfgBh}{0.65}   % half height of the 1/s block

% A summing node at (#1,#2).
\newcommand{\sfgSum}[2]{%
  \draw (#1,#2) circle (\sfgR);
  \draw ({#1-0.17},#2) -- ({#1+0.17},#2);
  \draw (#1,{#2-0.17}) -- (#1,{#2+0.17});
}

% A block centred at (#1,#2) carrying #3.
\newcommand{\sfgBox}[3]{%
  \draw ({#1-\sfgBw},{#2-\sfgBh}) rectangle ({#1+\sfgBw},{#2+\sfgBh});
  \node at (#1,#2) {#3};
}

\newcommand{\sfgDot}[2]{\fill (#1,#2) circle (0.07);}
